\section*{Control de cambios}
\phantomsection
\addcontentsline{toc}{section}{Control de cambios}
\noindent\textbf{Comparación:} \texttt{MANUAL\_DE\_COORDINACION\_ANTIGUO.pdf} frente a los fuentes de \texttt{DGP\_P1\_Bueno}.

Cada código tiene un color propio, aplicado también a los títulos y al contenido nuevo o modificado que se indica. Las figuras se identifican mediante un marco del mismo color, conservando la imagen original. Los cambios de presentación se describen en MC-01; el texto que no cambia mantiene su formato habitual. Las supresiones se documentan aquí porque ya no tienen un bloque que marcar.

\entradaCambio{MC-01}{cambioPresentacion}{Presentación y versión}
{La portada no indicaba versión ni fecha; el encabezado identificaba el manual mediante texto y varias secciones compartían página.}
{La portada incorpora un logotipo central, la indicación «VERSIÓN 2» y la fecha 05/11/2026. Se aplica la paleta CLENCH, se utiliza el logotipo en la cabecera y se añaden saltos de página antes de las secciones principales.}
{Portada, encabezados y títulos del documento; versión, fecha y marco del logotipo central marcados en gris pizarra.}

\entradaCambio{MC-02}{cambioArquitectura}{Diagrama de arquitectura}
{Recursos Software de Desarrollo terminaba con el ecosistema y las herramientas; no incluía un apartado de arquitectura.}
{Se añade «Diagrama de Arquitectura», con la descripción de cliente Flutter/Dart, servidor Python/FastAPI, persistencia PostgreSQL, diseño y apoyos visuales, y una figura de arquitectura.}
{Apartado 3.2, completo, y marco de su figura.}

\entradaCambio{MC-03}{cambioRevision}{Procedimiento de gestión de cambios}
{El plan se centraba en refactorizaciones y contenía tres pasos: notificación, análisis de impacto y decisión formal; admitía simplificar cambios insignificantes.}
{Se amplía a requisitos, alcance, diseño, código, accesibilidad, calendario y documentación. Se detallan seis pasos desde la solicitud hasta el cierre, las funciones de cada rol, las listas de verificación, los criterios de aceptación y la trazabilidad Jira--ramas--Pull Requests--versiones semánticas.}
{Sección 9 completa, incluidos los nuevos apartados 9.2 y 9.3.}

\entradaCambio{MC-04}{cambioRNF}{Plan de gestión de calidad}
{Se describían cinco pasos de calidad por iteración, un apartado de hardening y una tabla de herramientas (GitHub, SonarCloud/ESLint, Jest, Lighthouse/WAVE y Jira).}
{Se sustituye esa estructura por un plan con seis pasos, responsables, informe DC-03, registro de evidencias y estados de prueba, corrección y regresión, lista de verificación y criterios de aceptación. Se añaden pautas para distinguir actas internas y validaciones del centro y los datos mínimos de las actas. El apartado y la tabla anteriores dejan de figurar como bloques independientes.}
{Sección 11 completa; las pautas sobre actas procedían del contenido documental de la propuesta antigua.}

\entradaCambio{MC-05}{cambioInformacion}{Valoración del trabajo individual}
{Se indicaba únicamente el registro aproximado de horas en la plantilla compartida.}
{Se mantiene ese registro y se añade que se valorará positivamente el número de tareas realizadas por encima del tiempo registrado.}
{Párrafo de «Control del trabajo individual», apartado 12.1.}

\entradaCambio{MC-06}{cambioBeneficios}{Manual de accesibilidad y calidad incorporado}
{El documento terminaba tras la medición del desempeño y no incorporaba un manual formativo.}
{Se incorpora un manual con índice propio: propósito, diseño visual e interacción accesible, formación y buenas prácticas de backend, prácticas de calidad, formación, responsabilidades y lista de comprobación con referencias.}
{Sección 13 completa, incluidos los nueve bloques del manual integrado.}

\entradaCambio{MC-07}{cambioDocumentacion}{Actas del equipo y visita al cliente}
{No se incluían actas de reuniones en el manual antiguo.}
{Se incorporan ocho actas: visita al cliente del 02/10 y reuniones del 01/10, 09/10, 15/10, 19/10, 22/10, 29/10 y 02/11 de 2026, con participantes, desarrollo, acuerdos, compromisos y seguimiento.}
{Sección 14 completa, con sus tablas de asistencia y compromisos.}
